% =========================================================================
% MÓDULO 6: REFERENCIAS BIBLIOGRÁFICAS COMPLETAS (DIAPOSITIVAS 29 A 31)
% =========================================================================
\bibslidetrue
\setbeamertemplate{bibliography item}[text]
\begin{frame}[allowframebreaks]{Referencias Bibliogr\'aficas}{Fuentes cient\'ificas, metodol\'ogicas y de estado del arte citadas en la presentaci\'on}
  \scriptsize
  \begin{thebibliography}{21}
    \setlength{\itemsep}{3.5pt}

    \bibitem{miller2008prediction}
    K.~J. Miller and G.~Schalk, ``Prediction of finger flexion: 4th {Brain-Computer Interface} data competition ({Dataset 4}),'' \emph{BCI Competition IV}, 2008. [Online]. Available: \url{https://www.bbci.de/competition/iv/desc_4.pdf}

    \bibitem{schalk2004bci2000}
    G.~Schalk, D.~J. McFarland, T.~Hinterberger, N.~Birbaumer, and J.~R. Wolpaw, ``{BCI2000}: A general-purpose brain-computer interface ({BCI}) system,'' \emph{IEEE Transactions on Biomedical Engineering}, vol.~51, no.~6, pp. 1034--1043, 2004.

    \bibitem{Schalk2011}
    G.~Schalk and E.~C. Leuthardt, ``Brain-computer interfaces using electrocorticographic signals,'' \emph{IEEE Reviews in Biomedical Engineering}, vol.~4, pp. 140--154, 2011.

    \bibitem{kaplan2019decoding}
    K.~Volkova, M.~A. Lebedev, A.~Kaplan, and A.~Ossadtchi, ``Decoding movement from electrocorticographic activity: A review,'' \emph{Frontiers in Neuroinformatics}, vol.~13, p.~74, 2019.

    \bibitem{Bundy2016}
    D.~T. Bundy, M.~Pahwa, N.~Szrama, and E.~C. Leuthardt, ``Decoding three-dimensional reaching movements using electrocorticographic signals in humans,'' \emph{Journal of Neural Engineering}, vol.~13, no.~2, p. 026021, 2016.

    \bibitem{gruenwald2023highgamma}
    J.~Gruenwald, S.~Sieghartsleitner, C.~Kapeller, J.~Scharinger, K.~Kamada, P.~Brunner, and C.~Guger, ``Characterization of high-gamma activity in electrocorticographic signals,'' \emph{Frontiers in Neuroscience}, vol.~17, p. 1206120, 2023.

    \bibitem{wirth2000crispdm}
    R.~Wirth and J.~Hipp, ``{CRISP-DM}: Towards a standard process model for data mining,'' in \emph{Proc. 4th Int. Conf. Practical Applications of Knowledge Discovery and Data Mining}, vol.~1, 2000, pp. 29--39.

    \framebreak

    \bibitem{roberts2017crossvalidation}
    D.~R. Roberts \emph{et~al.}, ``Cross-validation strategies for data with temporal, spatial, hierarchical, or phylogenetic structure,'' \emph{Ecography}, vol.~40, no.~8, pp. 913--929, 2017.

    \bibitem{ronneberger2015unet}
    O.~Ronneberger, P.~Fischer, and T.~Brox, ``{U-Net}: Convolutional networks for biomedical image segmentation,'' in \emph{Medical Image Computing and Computer-Assisted Intervention ({MICCAI})}, vol. 9351, Springer, 2015, pp. 234--241.

    \bibitem{Sliwowski2022}
    M.~Sliwowski, ``Artificial intelligence for real-time decoding of motor commands from {ECoG} of disabled subjects for chronic brain computer interfacing,'' Ph.D. dissertation, Universit\'e Grenoble Alpes, 2022.

    \bibitem{loshchilov2019decoupled}
    I.~Loshchilov and F.~Hutter, ``Decoupled weight decay regularization,'' in \emph{International Conference on Learning Representations ({ICLR})}, 2019, pp. 1--19.

    \bibitem{Bergstra2011}
    J.~Bergstra, R.~Bardenet, Y.~Bengio, and B.~K\'egl, ``Algorithms for hyper-parameter optimization,'' in \emph{Advances in Neural Information Processing Systems ({NeurIPS})}, vol.~24, 2011, pp. 2546--2554.

    \bibitem{akiba2019optuna}
    T.~Akiba, S.~Sano, T.~Yanase, T.~Ohta, and M.~Koyama, ``{Optuna}: A next-generation hyperparameter optimization framework,'' in \emph{Proc. 25th ACM SIGKDD Int. Conf. Knowledge Discovery \& Data Mining}, 2019, pp. 2623--2631.

    \bibitem{cawley2010overfitting}
    G.~C. Cawley and N.~L.~C. Talbot, ``On over-fitting in model selection and subsequent selection bias in performance evaluation,'' \emph{Journal of Machine Learning Research}, vol.~11, pp. 2079--2107, 2010.

    \framebreak

    \bibitem{lin1989concordance}
    L.~I.-K. Lin, ``A concordance correlation coefficient to evaluate reproducibility,'' \emph{Biometrics}, vol.~45, no.~1, pp. 255--268, 1989.

    \bibitem{liang2012decoding}
    N.~Liang and L.~Bougrain, ``Decoding finger flexion from electrocorticographic signals using a sparse {G}aussian process,'' \emph{Frontiers in Neuroscience}, vol.~6, p.~91, 2012.

    \bibitem{Xie2018}
    Z.~Xie, O.~Schwartz, and A.~Prasad, ``Decoding of finger trajectory from {ECoG} using deep learning,'' \emph{Journal of Neural Engineering}, vol.~15, no.~3, p. 036009, 2018.

    \bibitem{Lomtev2022}
    V.~Lomtev, A.~Kovalev, and A.~Timchenko, ``{FingerFlex}: Inferring finger trajectories from {ECoG} signals,'' \emph{arXiv preprint arXiv:2211.01960}, 2022.

    \bibitem{wang2025dtcnet}
    F.~Wang, Z.~Luo, W.~Lv, and X.~Zhu, ``{DTCNet}: Finger flexion decoding with three-dimensional {ECoG} data,'' \emph{Frontiers in Computational Neuroscience}, vol.~19, p. 1627819, 2025.

    \bibitem{Xie2026}
    Y.~Xie \emph{et~al.}, ``{BrainDistill}: Implantable motor decoding with task-specific knowledge distillation,'' \emph{arXiv preprint arXiv:2601.17625}, 2026.

    \bibitem{rodas_valencia_tpe_ecog}
    S.~Rodas~Bautista and J.~Valencia~San~Roman, ``{TPE\_ECoG}: Optimizaci\'on mediante {TPE} para la decodificaci\'on de se\~nales {ECoG},'' GitHub, 2026. [Online]. Available: \url{https://github.com/SaulRodas/TPE_ECoG}

  \end{thebibliography}
\end{frame}
\bibslidefalse
